\makeabstract
    {
    Investigation of Student Experience with Computational Physics and Computational Astronomy Activities  
    }
    {
    Baokun Chen\textsuperscript{1}, Daniel Lu\textsuperscript{1}, Bahar Modir\textsuperscript{1}
    }
    {
    \textsuperscript{1} Department of Physics and Astronomy, Amherst College, Amherst, MA
    }
    {
    Modern day physics and astronomy research has increasingly relied on computational methods for solving complex problems. Accordingly, the Next Generation Science Standards has recognized computation as a priority for physics instruction. The Physics Education Research (PER) community has thought extensively about course design in introductory and upper-division physics, but computation in undergraduate research and specific practices have been understudied. We aim to examine students' experience with coding and programming in selected courses and undergraduate research. We ask: 1) How do physics/astronomy students from different levels of research and coding approach computational learning and problem solving? 2) How do introductory and self-explanatory materials can aid the overall learning experience of students in programming?
    }
    {
    End-of-Spring 2026 surveys were collected from 18 students: 10 from PHYS 125: Oscillations and Waves, two from ASTR 200: Intro to Data Science with Astronomical Applications, and six Physics and Astronomy thesis students. We conducted a 5-step thematic analysis using Notion and Dedoose: (1) reviewing a subset of transcripts, (2) summarizing responses, (3) generating initial codes, (4) applying and refining the coding scheme across the full dataset, and (5) grouping codes into overarching themes and inter-Rater Reliability was conducted to achieve a higher percentage of agreement.
    }
    {
    Through thematic analysis, we find 6 prominent themes across our data in students' experience with programming. Comparing beginner and research students (RQ1): We find that beginners see coding as less connected to their content learning/careers goals while research students see it as integral, though both groups value Python's data visualization capabilities. Research students also plan and structure their code more deliberately than beginners, and research students and beginner students have different debugging processes in terms of AI-usage and optimization. Examining introductory materials, we see (RQ2): current Python materials in PHYS 125 can be aided with more deliberate design for supporting students with various coding experience and background.
    }
    {
    Following our thematic analysis, we propose four recommendations for future course improvement to Incorporate large-scale projects, differentiate the intro file by skill level, provide TAs with preparation resources, and expand debugging-focused office hours. Possible limitations to our research include our small sample size and survey structure. In the future, we could refine the questions on the survey to better capture student experiences and expand to a larger sample size.
    }

    \newpage
    